\section{Problem Set 5: Vectors and Coordinates}

\begin{enumerate}
\item \textbf{Problem 1}\par
Explain the difference between the point \(P=(3,-2)\) and the vector \(v=(3,-2)\). What does each object describe?

\item \textbf{Problem 2}\par
For
\[
A=(1,4,-2),
\qquad
B=(5,-1,3),
\]
find \(\overrightarrow{AB}\) and \(\overrightarrow{BA}\).

\item \textbf{Problem 3}\par
Let
\[
u=(2,-1,3),
\qquad
v=(-4,5,1).
\]
Compute \(u+v\), \(u-v\), \(3u\), and \(2u-v\).

\item \textbf{Problem 4}\par
Find the position vector of the point \(P=(-2,5,1)\). Then find the endpoint of this vector if it begins at \(A=(3,-1,4)\) instead of the origin.

\item \textbf{Problem 5}\par
Find the length of
\[
v=(3,-4,12).
\]

\item \textbf{Problem 6}\par
Find the distance between
\[
A=(1,-2,3)
\quad\text{and}\quad
B=(5,1,-1).
\]

\item \textbf{Problem 7}\par
Find the midpoint of the segment joining
\[
A=(-3,4,2)
\quad\text{and}\quad
B=(5,-2,8).
\]

\item \textbf{Problem 8}\par
Find a unit vector in the direction of
\[
v=(2,-2,1).
\]

\item \textbf{Problem 9}\par
For
\[
u=(1,2,-1),
\qquad
v=(3,0,4),
\]
compute \(u\cdot v\).

\item \textbf{Problem 10}\par
Find the angle between
\[
u=(1,0,1),
\qquad
v=(1,1,0).
\]

\item \textbf{Problem 11}\par
Find the value of \(k\) for which
\[
u=(2,k,-1)
\quad\text{and}\quad
v=(1,3,5)
\]
are perpendicular.

\item \textbf{Problem 12}\par
Determine whether the triangle with vertices
\[
A=(0,0),\quad B=(4,0),\quad C=(4,3)
\]
is right-angled. Use vectors and the dot product.

\item \textbf{Problem 13}\par
Compute
\[
u\times v
\quad\text{for}\quad
u=(1,2,0),
\qquad
v=(3,-1,4).
\]

\item \textbf{Problem 14}\par
Verify that the vector found in the previous problem is perpendicular to both \(u\) and \(v\).

\item \textbf{Problem 15}\par
Find the area of the parallelogram spanned by
\[
u=(2,1,0),
\qquad
v=(1,-1,3).
\]
Then find the area of the triangle spanned by the same vectors.

\item \textbf{Problem 16}\par
Find a normal vector to the plane containing the direction vectors
\[
u=(1,2,-1),
\qquad
v=(2,0,3).
\]

\item \textbf{Problem 17}\par
Determine whether the points
\[
A=(1,2,3),\quad B=(3,6,5),\quad C=(5,10,7)
\]
are collinear.

\item \textbf{Problem 18}\par
A particle moves from \(A=(2,-1)\) by the displacement \(v=(5,3)\), and then by \(w=(-2,4)\). Find its final position and its total displacement from \(A\).

\item \textbf{Problem 19}\par
The midpoint of \(AB\) is \(M=(2,1,-3)\), and \(A=(-1,4,2)\). Find \(B\).

\item \textbf{Problem 20}\par
For the points
\[
A=(1,0,2),\quad B=(4,2,1),\quad C=(2,5,3),
\]
find \(\overrightarrow{AB}\) and \(\overrightarrow{AC}\), compute their dot product and cross product, determine the angle at \(A\), and find the area of triangle \(ABC\).
\end{enumerate}

\section{Problem Set 6: Lines}

\begin{enumerate}
\item \textbf{Problem 1}\par
Write a vector-parametric equation of the line through
\[
A=(2,-1)
\]
with direction vector
\[
v=(3,4).
\]

\item \textbf{Problem 2}\par
Find a parametric equation of the line through
\[
A=(-1,2)
\quad\text{and}\quad
B=(5,-4).
\]

\item \textbf{Problem 3}\par
Determine whether the point \(P=(7,3)\) lies on the line
\[
(x,y)=(1,-1)+t(2,1).
\]

\item \textbf{Problem 4}\par
Write the coordinate parametric equations of
\[
P=(2,3)+t(-1,4).
\]
Find the points corresponding to \(t=-2\), \(t=0\), and \(t=3\).

\item \textbf{Problem 5}\par
Find the point-slope equation of the line through \((3,-2)\) with slope \(5\).

\item \textbf{Problem 6}\par
Convert
\[
y-4=-2(x+1)
\]
to slope-intercept form and general form.

\item \textbf{Problem 7}\par
Write equations for:
\begin{enumerate}
\item the vertical line through \((3,-1)\);
\item the horizontal line through \((-2,5)\).
\end{enumerate}

\item \textbf{Problem 8}\par
Convert the line
\[
(x,y)=(1,2)+t(3,-1)
\]
to general form.

\item \textbf{Problem 9}\par
Convert
\[
2x-5y+10=0
\]
to a parametric equation.

\item \textbf{Problem 10}\par
Determine whether the lines
\[
\ell_1:(x,y)=(1,0)+t(2,3),
\qquad
\ell_2:(x,y)=(-1,4)+s(4,6)
\]
are distinct parallel lines or the same line.

\item \textbf{Problem 11}\par
Determine whether
\[
\ell_1:2x-y+3=0
\quad\text{and}\quad
\ell_2:x+2y-4=0
\]
are perpendicular.

\item \textbf{Problem 12}\par
Find the intersection point of
\[
3x-y-5=0,
\qquad
x+2y-8=0.
\]

\item \textbf{Problem 13}\par
Find the distance from \(P=(4,-1)\) to the line
\[
2x+3y-6=0.
\]

\item \textbf{Problem 14}\par
For the oriented line
\[
x-2y+1=0,
\]
determine on which side of the line the points \(P=(0,0)\) and \(Q=(4,1)\) lie. Are they on the same side?

\item \textbf{Problem 15}\par
Write a parametric equation of the line in space through
\[
A=(1,-2,3)
\]
with direction vector
\[
v=(2,0,-1).
\]

\item \textbf{Problem 16}\par
Convert
\[
(x,y,z)=(2,1,-3)+t(4,-2,5)
\]
to symmetric form.

\item \textbf{Problem 17}\par
Write an appropriate symmetric description of
\[
(x,y,z)=(1,-2,3)+t(4,0,-1),
\]
taking proper account of the zero component.

\item \textbf{Problem 18}\par
Classify the following pair of lines in space as parallel, intersecting, or skew:
\[
\ell_1:(x,y,z)=(0,0,0)+t(1,1,0),
\]
\[
\ell_2:(x,y,z)=(1,0,1)+s(0,1,1).
\]

\item \textbf{Problem 19}\par
A line in the plane is given by a point \(A=(2,1)\) and a normal vector \(n=(3,-4)\). Write its normal equation, general equation, and one direction vector.

\item \textbf{Problem 20}\par
The line passes through \(A=(-1,3)\) and \(B=(5,-1)\). Find:
\begin{enumerate}
\item a direction vector;
\item a parametric equation;
\item a normal vector;
\item a general equation;
\item its slope, if defined;
\item the distance from the origin to the line.
\end{enumerate}
\end{enumerate}

\section{Problem Set 7: Planes}

\begin{enumerate}
\item \textbf{Problem 1}\par
Write the normal equation and general equation of the plane through
\[
P_0=(1,-2,3)
\]
with normal vector
\[
n=(2,1,-1).
\]

\item \textbf{Problem 2}\par
For the plane
\[
3x-2y+4z-7=0,
\]
state a normal vector and find one convenient point on the plane.

\item \textbf{Problem 3}\par
Determine whether \(P=(2,1,-1)\) lies on
\[
x+2y-z-5=0.
\]

\item \textbf{Problem 4}\par
Write a parametric equation of the plane through
\[
P_0=(1,0,2)
\]
with direction vectors
\[
u=(1,2,0),
\qquad
v=(0,1,3).
\]

\item \textbf{Problem 5}\par
Find a parametric equation of the plane through
\[
A=(1,0,0),\quad B=(0,2,0),\quad C=(0,0,3).
\]

\item \textbf{Problem 6}\par
Find a normal vector to the plane through
\[
A=(1,0,1),\quad B=(3,1,0),\quad C=(0,2,4).
\]

\item \textbf{Problem 7}\par
Find the general equation of the plane through the three points in the previous problem.

\item \textbf{Problem 8}\par
Convert the parametric plane
\[
P=(1,2,0)+s(1,0,1)+t(0,2,-1)
\]
to general form.

\item \textbf{Problem 9}\par
Convert
\[
2x-y+z-4=0
\]
to one possible parametric form.

\item \textbf{Problem 10}\par
Determine whether the planes
\[
\pi_1:x-2y+3z-1=0,
\qquad
\pi_2:2x-4y+6z+5=0
\]
are distinct parallel planes or the same plane.

\item \textbf{Problem 11}\par
Determine whether
\[
\pi_1:x+y-z=0,
\qquad
\pi_2:2x-y+z-3=0
\]
are perpendicular.

\item \textbf{Problem 12}\par
Find the acute angle between
\[
\pi_1:x+2y+2z=0
\quad\text{and}\quad
\pi_2:2x-y+2z-4=0.
\]

\item \textbf{Problem 13}\par
Find the distance from
\[
P=(3,-1,2)
\]
to the plane
\[
2x-y+2z-7=0.
\]

\item \textbf{Problem 14}\par
Find the intersection point of the line
\[
(x,y,z)=(1,0,2)+t(2,-1,1)
\]
with the plane
\[
x+y+z-6=0.
\]

\item \textbf{Problem 15}\par
Determine whether the line
\[
(x,y,z)=(1,2,0)+t(1,-1,2)
\]
lies in, is parallel to, or intersects the plane
\[
2x+2y-z-6=0.
\]

\item \textbf{Problem 16}\par
Find the equation of the plane through \(P=(2,-1,4)\) parallel to
\[
3x-y+2z+5=0.
\]

\item \textbf{Problem 17}\par
Find a plane through \(P=(1,2,3)\) that is perpendicular to the plane
\[
x-y+z=0.
\]
Give one valid answer and explain why it works.

\item \textbf{Problem 18}\par
Find a parametric equation of the line of intersection of
\[
x+y+z=3,
\qquad
x-y+z=1.
\]

\item \textbf{Problem 19}\par
A plane contains the line
\[
\ell:(x,y,z)=(1,0,2)+t(2,1,-1)
\]
and the point \(Q=(0,3,1)\). Find a general equation of the plane.

\item \textbf{Problem 20}\par
For
\[
A=(1,0,0),\quad B=(0,1,0),\quad C=(0,0,1),
\]
find a parametric and a general equation of the plane through the points, verify that all three points satisfy the equation, find the distance from the origin to the plane, and determine where the line
\[
(x,y,z)=t(1,1,1)
\]
intersects it.
\end{enumerate}
